%!TEX TS-program = xelatex
\documentclass[12pt,a4paper]{article}
\usepackage{ctex}
\usepackage{amsmath,amsfonts,amssymb}
\usepackage{graphicx}
\usepackage{booktabs} % 三线表
\usepackage{multirow} % 合并单元格
\usepackage{xcolor}   % 颜色
\usepackage{hyperref} % 链接

\title{\LaTeX{} 基础排版综合练习}
\author{你的名字}
\date{\today}

\begin{document}
\maketitle

\section{基础文本排版}
熟悉不同的字体样式是 \LaTeX{} 的基本功。
这里是 \textbf{粗体}，这里是 \textit{斜体}，这里是 \underline{下划线}。
当你需要强调时，可以使用 \emph{强调文字}。
文字颜色可以修改为 \textcolor{red}{红色} 或 \textcolor{blue}{蓝色}。

换段落时，只需要在源代码中\textbf{空一行}即可。

这是新的一段。如果要强制换行但是不换段，可以使用双反斜杠 \\
这是由于双反斜杠造成的换行。

\section{数学符号与公式}
数学公式分为两种：
\begin{itemize}
    \item \textbf{行内公式：} 直接嵌在段落里，比如著名的欧拉公式 $e^{i\pi} + 1 = 0$。
    \item \textbf{独立公式（行间公式）：} 会单独占一行并居中，通常带编号：
\end{itemize}

\begin{equation}
    f(x) = \int_{-\infty}^{\infty} \hat{f}(\xi)\,e^{2 \pi i \xi x} \,d\xi
\end{equation}

还有常用的多行对齐公式：
\begin{align}
    (x+y)^3 &= (x+y)(x^2 + 2xy + y^2) \nonumber \\ 
            &= x^3 + 3x^2y + 3xy^2 + y^3
\end{align}

\section{列表项练习}
这是有序列表的支持：
\begin{enumerate}
    \item 掌握文本排版
    \item 掌握数学公式
    \item 掌握插图和表格
\end{enumerate}

\section{表格排版}
在学术论文中最常用的是“三线表”。可以使用 \texttt{booktabs} 宏包实现：

\begin{table}[htbp]
    \centering
    \caption{常见的机器学习模型性能对比}
    \label{tab:models}
    \begin{tabular}{lccc}
        \toprule
        \textbf{模型名称} & \textbf{准确率} & \textbf{召回率} & \textbf{F1 分数} \\
        \midrule
        逻辑回归 & 85.2\% & 83.4\% & 84.3\% \\
        支持向量机 (SVM) & 88.7\% & 87.1\% & 87.9\% \\
        随机森林 & 92.4\% & 90.8\% & 91.6\% \\
        \bottomrule
    \end{tabular}
\end{table}

\section{插入图片}
\LaTeX{} 插入图片通常使用 \texttt{graphicx} 宏包。被注释掉的代码展示了插入图片的典型方法：

% \begin{figure}[htbp]
%     \centering
%     \includegraphics[width=0.6\linewidth]{image_name.png} % 将此处的文件名替换为真实图片名
%     \caption{这是一张插入的示例图片}
%     \label{fig:example}
% \end{figure}

如图 \ref{tab:models} 所示（这里演示了交叉引用功能），你可以轻松引用文档中的表格或图片。

\end{document}
